\documentclass[11pt]{article}

%------------------------------------------------
% Packages
%------------------------------------------------
\usepackage[a4paper,margin=2.5cm]{geometry}
\usepackage{amsmath,amssymb,amsthm}
\usepackage{mathtools}
\usepackage{enumitem}
\usepackage{hyperref}

\hypersetup{
    colorlinks=true,
    linkcolor=blue,
    citecolor=blue,
    urlcolor=blue
}

%------------------------------------------------
% Title Information
%------------------------------------------------

\title{\bfseries Calmness of the Solution-Set Mapping\\for Linear
  Bilevel and Pricing Problems}

%%%%%%%%%%%%%%%%%%%%%%%%%%%%%%%%%%%%%%%%%%%%%%%%%%%%%%%%%%%%%%%%%%%%%%%%%%%

\author{
  Henri Lefebvre\thanks{LIRMM, University of Montpellier, CNRS,
    Montpellier, France,
    \texttt{henri.lefebvre@cnrs.fr}}
  \and
  Martin Schmidt\thanks{Trier University,
    Department of Mathematics,
    Universitätsring 15,
    54296 Trier,
    Germany,
    \texttt{martin.schmidt@uni-trier.de}}
  \and
  Anton Svensson\thanks{Universidad de O'Higgins, Av Libertador Bernardo O'Higgins 611,
    Rancagua,
    Chile,
    \texttt{anton.svensson@uoh.cl}}
  \and
  Johannes Thürauf\thanks{University of Technology Nuremberg (UTN),
    Department Liberal Arts and Social Sciences,
    Discrete Optimization,
    Dr.-Luise-Herzberg-Str.\ 4,
    90461 Nuremberg,
    Germany,
    \texttt{johannes.thuerauf@utn.de}}
}

\date{}

%------------------------------------------------

\begin{document}

\maketitle

\vspace{-1em}

\noindent
\textbf{Keywords:}
Bilevel optimization,
Linear bilevel optimization,
Pricing problems,
Approximate feasibility,
Calmness,
Outer semicontinuity

\vspace{0.5em}


\vspace{1em}

%%%%%%%%%%%%%%%%%%%%%%%%%%%%%%%%%%%%%%%%%%%%%%%%%%
\section*{Abstract}

We study linear bilevel and pricing problems in which the upper- and
lower-level constraints' right-hand sides are perturbed.
In this setting, it is an important question, also for the validity of
numerical solution schemes, if the solution-set mapping of the
parametric bilevel problem is calm at the zero-perturbation.
We provide the complete picture both for linear bilevel as well as for
pricing problems.
If the result is positive or not depends on whether the problems have
coupling constraints or not and on whether the perturbation is allowed
to take place in both levels and if they lead to relaxations or
tightenings of the respective constraints.
In particular, the solution-set mapping is calm for linear bilevel
problems without coupling constraints and for pricing problems if the
upper-level problem is not tightened by the perturbation.
For the negative results, we provide illustrative counterexamples.

%%%%%%%%%%%%%%%%%%%%%%%%%%%%%%%%%%%%%%%%%%%%%%%%%%
\vspace{1em}

\noindent
\textbf{Acknowledgements.}
We acknowledge the support by the German
Bundesministerium f\"ur Bildung und Forschung within
the project ``RODES'' (F\"orderkennzeichen 05M22UTB).
Moreover, the authors thank the DFG for their support within
the research training group~2126 ``Algorithmic Optimization''.
The third author acknowledges the support of ANID-Chile through the
projects Fondecyt regular 1251159 and Fondecyt iniciación 11261583.

%%%%%%%%%%%%%%%%%%%%%%%%%%%%%%%%%%%%%%%%%%%%%%%%%%
\vspace{1em}

%\begin{thebibliography}{9}
%
%\bibitem{Attouch2013}
%H.~Attouch, J.~Bolte and B.~Svaiter.
%\newblock Convergence of descent methods for semi-algebraic and tame problems.
%\newblock {\em Mathematical Programming}, 137 (2013), 91--129.
%
%\bibitem{Nesterov2018}
%Y.~Nesterov.
%\newblock {\em Lectures on Convex Optimization}.
%\newblock Springer, 2018.
%
%\end{thebibliography}

\end{document}